\section{An all-length Rauzy matrix decoder}
\label{sec:decoder}

The main operator proof can isolate a branch by support and therefore does
not require matrix injectivity.  The next theorem nevertheless closes a
second proposed alternative: different chronological branches cannot project to
the same full matrix and cancel through opposite metaplectic signs in the
present full-rank four-letter model.

The geometric antecedent is the standard simplicial-cylinder coding of
Rauzy induction: finite words determine nested projective cones, with the
two first-level cone interiors separated by the comparison hyperplane
\citep{kerckhoff1985}.  We do not claim novelty for that coding geometry.
The statement below is a self-contained algorithmic form of it, tailored to
the later-on-the-left matrix convention: a single cumulative matrix is
inverted by a deterministic row-dominance subtraction.

\begin{theorem}[Fixed-start path decoding]
\label{thm:decoder}
Fix an irreducible labeled Rauzy permutation \(\pi\).  For every finite path
\(\gamma\) starting at \(\pi\), the pair \((\pi,B_\gamma)\), with the
conventions \eqref{eq:edge-matrix}--\eqref{eq:chronology}, uniquely
determines every directed arrow of \(\gamma\).
\end{theorem}

\begin{proof}
Let \(\gamma=e_1\cdots e_n\) and put \(R=B_\gamma^T\).  If the first edge
has winner \(w\) and loser \(\ell\), then
\begin{equation}
 R=(I+E_{w,\ell})R',
 \qquad R'=B_{e_2\cdots e_n}^T.
 \label{eq:first-factor}
\end{equation}
Consequently
\begin{equation}
 \operatorname{row}_w(R)
 =\operatorname{row}_w(R')+\operatorname{row}_\ell(R'),
 \qquad
 \operatorname{row}_\ell(R)=\operatorname{row}_\ell(R').
 \label{eq:row-dominance}
\end{equation}
Every path matrix is nonnegative and unimodular.  The true first edge
therefore satisfies componentwise winner-row dominance
\begin{equation}
 \operatorname{row}_w(R)\ge\operatorname{row}_\ell(R),
 \label{eq:dominance}
\end{equation}
with at least one strict coordinate.

At a fixed irreducible state, the two outgoing arrows exchange the same
rightmost top and bottom labels, so the other candidate requires the reverse
dominance.  If both candidates passed, the two rows would be equal.  That is
impossible because \(R\) is invertible.  Thus exactly one candidate identifies
the true first edge.

Peeling that edge is exact:
\begin{equation}
 R'=(I-E_{w,\ell})R.
 \label{eq:row-peel}
\end{equation}
The operation replaces the winner row by winner minus loser, leaves a
nonnegative integral matrix, and decreases the sum of all entries by the
strictly positive loser-row sum.  We update the labeled permutation and
repeat.  The entry sum cannot decrease indefinitely, and a valid path ends
at the identity.  Induction recovers the unique full word.
\end{proof}

The proof yields a reproducible decoder.

\begin{quote}
\textbf{Decoder.} Start with the labeled state \(\pi\) and
\(R=B_\gamma^T\).  Until \(R=I\): (i) test the two outgoing winner/loser
row dominances; (ii) select the unique passing edge; (iii) replace its winner
row by winner minus loser; (iv) advance the labeled state.  Reject an input
matrix if neither or both candidates pass, the subtraction becomes negative,
or the terminal matrix is not \(I\).
\end{quote}

\begin{corollary}[Return-matrix freeness]
\label{cor:freeness}
Return paths based at a fixed labeled state have distinct full Rauzy
matrices.  Concatenations of first-return paths form a free monoid on the
countable first-return generators.
\end{corollary}

\begin{proof}
Theorem~\ref{thm:decoder} recovers the complete edge word.  Visits to the
fixed state then give its unique decomposition into first returns.
\end{proof}

The word ``full'' matters.  A characteristic polynomial, spectrum, or
conjugacy class does not carry the row data required by
\eqref{eq:dominance}.  In a general stratum, restriction from relative to
absolute homology can also discard information.  For the present
four-letter \(\mathcal H(2)\) model, \(\det\Omega_\pi=1\), so no relative
kernel is present.  Matrix decoding recovers the path; the central
metaplectic sign is then recovered by multiplying the declared edge lifts.
We do not claim that an endpoint matrix without this lift convention already
contains that sign.

The decoder also explains why retaining genuine chronology is practical.
The ordered integer product is not merely a coarse endpoint statistic in
this fixed-start coding: it admits a deterministic inverse.  Replacing it by
an averaged transition matrix, a spectrum, or a Parikh count destroys that
inverse.
